\chapter{Integração}\label{ch:g12:integ}

A integração responde a duas perguntas ao mesmo tempo: qual é a área sob
uma curva e como se recupera uma \omterm{def:g11:func:function}{função} a partir da sua taxa de variação?
O teorema fundamental do cálculo afirma que essas são a mesma pergunta ---
a descoberta mais profunda e mais útil da matemática do século XVII.

\section{A integral de uma função contínua}

\begin{definition}[Integral de uma função positiva]\label{def:g12:integ:area}
Seja $f$ \omterm{def:g12:limcont:continuity}{contínua} e positiva em $\intcc{a}{b}$. A
\emph{integral}\index{integral}
\[
\int_a^b f(x)\,\dd x
\]
é a área, em unidades de área, da região delimitada pela curva de $f$,
pelo eixo $x$ e pelas retas verticais $x = a$ e $x = b$.
\end{definition}

\begin{omfigure}
\begin{tikzpicture}
\begin{axis}[
  width=8.4cm, height=5cm,
  axis lines=middle, xmin=-0.4, xmax=3.4, ymin=-0.3, ymax=2.3,
  xtick={0.5,2.7}, xticklabels={$a$,$b$}, ytick=\empty,
  samples=100, domain=-0.2:3.3]
  \addplot[name path=curve, omDef, thick, domain=0.2:3.2]
    {1 + 0.5*sin(deg(2*x)) + 0.2*x};
  \path[name path=axis] (0.5,0) -- (2.7,0);
  \addplot[omDef!20] fill between[of=curve and axis, soft clip={domain=0.5:2.7}];
  \node at (axis cs:1.6,0.55) {$\displaystyle\int_a^b f$};
\end{axis}
\end{tikzpicture}
\end{omfigure}

Para uma \omterm{def:g11:func:function}{função} de sinal qualquer, as áreas abaixo do eixo $x$ são contadas
negativamente; e define-se
$\int_b^a f(x)\,\dd x = -\int_a^b f(x)\,\dd x$.

\begin{omfigure}
\begin{tikzpicture}
\begin{axis}[omaxis,
  width=9.6cm, height=4.8cm,
  xmin=-0.5, xmax=7, ymin=-1.4, ymax=1.4,
  xtick={3.1416,6.2832}, xticklabels={$\pi$,$2\pi$}, ytick={-1,1}]
  \addplot[name path=sine, omDef, thick, domain=0:6.2832, samples=120]
    {sin(deg(x))};
  \path[name path=xaxis] (0,0) -- (6.2832,0);
  \addplot[omDef!20] fill between[of=sine and xaxis,
    soft clip={domain=0:3.1416}];
  \addplot[omThm!20] fill between[of=sine and xaxis,
    soft clip={domain=3.1416:6.2832}];
  \node at (axis cs:1.57,0.4) {$+$};
  \node at (axis cs:4.71,-0.4) {$-$};
\end{axis}
\end{tikzpicture}

{\small Áreas com sinal: $\int_0^{2\pi} \sin x\,\dd x = 0$, com a região
abaixo do eixo (vermelha) cancelando a região acima (azul).}
\end{omfigure}

\begin{proposition}[Propriedades da integral]\label{prop:g12:integ:properties}
Sejam $f, g$ \omterm{def:g12:limcont:continuity}{contínuas} em um \omterm{def:g10:numbers:interval}{intervalo} que contém $a, b, c$, e seja
$\lambda \in \R$.
\begin{enumerate}
  \item \emph{Linearidade:}
        $\displaystyle\int_a^b (f + \lambda g) = \int_a^b f + \lambda \int_a^b g$.
  \item \emph{Relação de Chasles:}
        $\displaystyle\int_a^c f = \int_a^b f + \int_b^c f$.
  \item \emph{Positividade:} se $a \leq b$ e $f \geq 0$ em
        $\intcc{a}{b}$, então $\displaystyle\int_a^b f \geq 0$; se
        $f \leq g$, então $\displaystyle\int_a^b f \leq \int_a^b g$.
\end{enumerate}
\end{proposition}

\begin{proof}
Chasles e a positividade são imediatas pela interpretação como área (as
áreas se somam quando as regiões são justapostas; uma região de altura
positiva tem área positiva). A comparação segue aplicando a positividade a
$g - f$. A linearidade é intuitivamente clara para a soma de funções
positivas (empilhe as áreas) e é demonstrada com rigor na graduação;
\emph{ela é admitida aqui}.
\end{proof}

\section{O teorema fundamental do cálculo}

\begin{definition}[Primitiva]\label{def:g12:integ:primitive}
Uma \emph{primitiva}\index{primitiva} (ou antiderivada) de $f$ em um
\omterm{def:g10:numbers:interval}{intervalo} $I$ é uma \omterm{def:g11:func:function}{função} \omterm{def:g12:deriv:derivative}{derivável} $F$ em $I$ tal que $F' = f$.
\end{definition}

\begin{proposition}\label{prop:g12:integ:primunique}
Se $F$ é uma \omterm{def:g12:integ:primitive}{primitiva} de $f$ em um \omterm{def:g10:numbers:interval}{intervalo} $I$, as \omterm{def:g12:integ:primitive}{primitivas} de $f$ em
$I$ são exatamente as funções $F + c$, $c \in \R$. Dados $x_0 \in I$ e
$y_0 \in \R$, existe uma única \omterm{def:g12:integ:primitive}{primitiva} com $F(x_0) = y_0$.
\end{proposition}

\begin{proof}
Se $G' = F' = f$, então $(G - F)' = 0$ no \omterm{def:g10:numbers:interval}{intervalo} $I$, de modo que
$G - F$ é constante.\footnote{Que uma \omterm{def:g11:func:function}{função} de \omterm{def:g12:deriv:derivative}{derivada} nula em um
\omterm{def:g10:numbers:interval}{intervalo} seja constante segue do \cref{thm:g12:deriv:variations}: ela é ao
mesmo tempo \omterm{def:g12:seq:monotonic}{crescente} e \omterm{def:g10:functions:variations}{decrescente}.} Reciprocamente, toda $F + c$ é
\omterm{def:g12:integ:primitive}{primitiva}. A condição $F(x_0) = y_0$ fixa a constante.
\end{proof}

\begin{theorem}[Teorema fundamental do cálculo]\label{thm:g12:integ:ftc}
\index{teorema fundamental do cálculo}
Seja $f$ \omterm{def:g12:limcont:continuity}{contínua} em um \omterm{def:g10:numbers:interval}{intervalo} $I$ e seja $a \in I$. A \omterm{def:g11:func:function}{função}
\[
F \colon x \longmapsto \int_a^x f(t)\,\dd t
\]
é a \omterm{def:g12:integ:primitive}{primitiva} de $f$ em $I$ que se anula em $a$. Por consequência, para
qualquer \omterm{def:g12:integ:primitive}{primitiva} $G$ de $f$ e $a, b \in I$:
\[
\int_a^b f(t)\,\dd t = \bigl[G(t)\bigr]_a^b = G(b) - G(a).
\]
\end{theorem}

\begin{omfigure}
\begin{tikzpicture}
\begin{axis}[omaxis,
  width=9.2cm, height=5.6cm,
  xmin=-0.4, xmax=4.2, ymin=-0.4, ymax=2.9,
  xtick={0.5,2.2,2.7}, xticklabels={$a$,$x$,\ $x+h$}, ytick=\empty]
  \addplot[name path=curve, omDef, thick, domain=0:3.9, samples=80]
    {1 + 0.4*x + 0.2*sin(deg(2*x))};
  \path[name path=xaxis] (0,0) -- (3.9,0);
  \addplot[omDef!20] fill between[of=curve and xaxis,
    soft clip={domain=0.5:2.2}];
  \addplot[omProp!40] fill between[of=curve and xaxis,
    soft clip={domain=2.2:2.7}];
  \draw[black, dashed, thin] (axis cs:2.2,1.69) -- (axis cs:2.7,1.69);
  \draw[black, dashed, thin] (axis cs:2.2,1.925) -- (axis cs:2.7,1.925)
    node[above right, font=\small] {$f(x+h)$};
  \node at (axis cs:1.35,0.6) {$F(x)$};
\end{axis}
\end{tikzpicture}

{\small A ideia da demonstração: o acréscimo $F(x+h) - F(x)$ é a área da
faixa estreita (laranja), espremida entre os retângulos de alturas $f(x)$
e $f(x+h)$ sobre uma base de comprimento $h$.}
\end{omfigure}

\begin{proof}[Demonstração no caso em que $f$ é crescente]
Fixe $x \in I$ e $h > 0$ com $x + h \in I$. Por Chasles,
\[
F(x+h) - F(x) = \int_x^{x+h} f(t)\,\dd t .
\]
Como $f$ é \omterm{def:g12:seq:monotonic}{crescente}, $f(x) \leq f(t) \leq f(x+h)$ para
$t \in \intcc{x}{x+h}$ e, por comparação de \omterm{def:g12:integ:area}{integrais} (a \omterm{def:g12:integ:area}{integral} de uma
constante $c$ em um \omterm{def:g10:numbers:interval}{intervalo} de comprimento $h$ vale $ch$):
\[
h\,f(x) \leq F(x+h) - F(x) \leq h\,f(x+h),
\]
de modo que
\[
f(x) \leq \frac{F(x+h) - F(x)}{h} \leq f(x+h).
\]
Quando $h \to 0^+$, $f(x+h) \to f(x)$ por \omterm{def:g12:limcont:continuity}{continuidade}, e o teorema do
confronto dá que o quociente de diferenças tende a $f(x)$; o caso $h < 0$ é
simétrico. Logo $F' = f$, e $F(a) = 0$. O caso \omterm{def:g12:limcont:continuity}{contínuo} geral (não
\omterm{def:g12:seq:monotonic}{monótono}) é demonstrado na graduação.

Por fim, se $G$ é uma \omterm{def:g12:integ:primitive}{primitiva} qualquer, $G = F + c$
(\cref{prop:g12:integ:primunique}), de modo que
$G(b) - G(a) = F(b) - F(a) = \int_a^b f$.
\end{proof}

\begin{example}
$\displaystyle\int_0^1 x^2\,\dd x = \left[\frac{x^3}{3}\right]_0^1 =
\frac13$: a área sob a \omterm{def:g11:quad:parabola}{parábola} é um terço do quadrado unitário, como
Arquimedes sabia.
\end{example}

A tabela de \omterm{def:g12:integ:primitive}{primitivas} se lê na tabela de \omterm{def:g12:deriv:derivative}{derivadas} ($u$ denota uma \omterm{def:g11:func:function}{função}
\omterm{def:g12:deriv:derivative}{derivável} e $c$ uma constante arbitrária):
\begin{center}
\renewcommand{\arraystretch}{1.7}
\begin{tabular}{c|c||c|c}
$f(x)$ & \omterm{def:g12:integ:primitive}{primitiva} & $f$ & \omterm{def:g12:integ:primitive}{primitiva} \\
\hline
$x^n \ (n \neq -1)$ & $\dfrac{x^{n+1}}{n+1} + c$ &
$u'u^n \ (n \neq -1)$ & $\dfrac{u^{n+1}}{n+1} + c$\\
$\dfrac1x \ (x > 0)$ & $\ln x + c$ &
$\dfrac{u'}{u} \ (u > 0)$ & $\ln u + c$\\
$\eu^x$ & $\eu^x + c$ &
$u'\eu^u$ & $\eu^u + c$\\
$\cos x$ & $\sin x + c$ &
$\dfrac{u'}{\sqrt u} \ (u>0)$ & $2\sqrt u + c$\\
$\sin x$ & $-\cos x + c$ & &
\end{tabular}
\end{center}

\begin{method}[Reconhecer a forma $u' \times (\text{algo em } u)$]\label{met:g12:integ:uprime}
Para integrar um produto, procure um fator que seja a \omterm{def:g12:deriv:derivative}{derivada} de uma
\omterm{def:g11:func:function}{função} interna $u$, a menos de uma constante multiplicativa. Por exemplo,
em $\int_0^1 x\,\eu^{x^2}\dd x$, o fator $x$ é $\frac12 (x^2)'$:
\[
\int_0^1 x\,\eu^{x^2}\dd x = \frac12\left[\eu^{x^2}\right]_0^1
= \frac{\eu - 1}{2}.
\]
\end{method}

\begin{theorem}[Integração por partes]\label{thm:g12:integ:ibp}
\index{integração por partes}
Sejam $u, v$ deriváveis em $\intcc{a}{b}$ com \omterm{def:g12:deriv:derivative}{derivadas} \omterm{def:g12:limcont:continuity}{contínuas}. Então
\[
\int_a^b u'(t)\,v(t)\,\dd t
= \bigl[u(t)\,v(t)\bigr]_a^b - \int_a^b u(t)\,v'(t)\,\dd t .
\]
\end{theorem}

\begin{proof}
A regra do produto dá $(uv)' = u'v + uv'$; integrar os dois lados em
$\intcc{a}{b}$ e usar o teorema fundamental no lado esquerdo dá
$\bigl[uv\bigr]_a^b = \int_a^b u'v + \int_a^b uv'$.
\end{proof}

\begin{example}\label{ex:g12:integ:ibp}
$\displaystyle\int_0^1 t\,\eu^{t}\,\dd t$: tome $u' = \eu^t$, $v = t$,
logo $u = \eu^t$, $v' = 1$:
\[
\int_0^1 t\,\eu^t\,\dd t = \bigl[t\,\eu^t\bigr]_0^1 - \int_0^1 \eu^t\,\dd t
= \eu - (\eu - 1) = 1 .
\]
\end{example}

\section{Aplicações}

\begin{definition}[Valor médio]\label{def:g12:integ:mean}
O \emph{valor médio}\index{valor médio} de uma \omterm{def:g11:func:function}{função} \omterm{def:g12:limcont:continuity}{contínua} $f$ em
$\intcc{a}{b}$ ($a<b$) é
\[
\mu = \frac{1}{b-a}\int_a^b f(t)\,\dd t .
\]
\end{definition}

\begin{proposition}\label{prop:g12:integ:meanbounds}
Se $m \leq f \leq M$ em $\intcc{a}{b}$, então $m \leq \mu \leq M$.
\end{proposition}

\begin{proof}
Integre as desigualdades $m \leq f(t) \leq M$ em $\intcc{a}{b}$ e divida
por $b - a > 0$.
\end{proof}

\begin{method}[Área entre duas curvas]\label{met:g12:integ:areabetween}
Se $f \geq g$ em $\intcc{a}{b}$, a área entre as duas curvas é
$\int_a^b \bigl(f(x) - g(x)\bigr)\dd x$. Se as curvas se cruzam, divida o
\omterm{def:g10:numbers:interval}{intervalo} nos pontos de cruzamento e integre $\abs{f - g}$ pedaço a
pedaço.
\end{method}

\begin{omfigure}
\begin{tikzpicture}
\begin{axis}[omaxis,
  width=8.4cm, height=5.8cm,
  xmin=-2.1, xmax=2.9, ymin=-1.8, ymax=5.6,
  xtick={-1,2}, ytick=\empty]
  \addplot[name path=line, omThm, thick, domain=-1.7:2.5] {x + 1};
  \addplot[name path=parab, omDef, thick, domain=-1.9:2.45] {x^2 - 1};
  \addplot[omDef!20] fill between[of=line and parab,
    soft clip={domain=-1:2}];
  \fill (axis cs:-1,0) circle (1.5pt);
  \fill (axis cs:2,3) circle (1.5pt);
\end{axis}
\end{tikzpicture}

{\small A área entre a reta $y = x + 1$ (vermelha) e a \omterm{def:g11:quad:parabola}{parábola}
$y = x^2 - 1$ (azul) é $\int_{-1}^{2}
\bigl((x+1) - (x^2-1)\bigr)\dd x = \frac92$.}
\end{omfigure}

\section{Exercícios}

\begin{exercise}[$\star$]\label{exo:g12:integ:1}
Calcule
\[
\int_1^2 \left(3x^2 - \frac{1}{x^2}\right)\dd x, \qquad
\int_0^{\pi/2} \cos t \,\dd t, \qquad
\int_0^{1} \frac{\dd t}{2t+1} .
\]
\end{exercise}

\begin{exercise}[$\star$]\label{exo:g12:integ:2}
Encontre a \omterm{def:g12:integ:primitive}{primitiva} $F$ de $f(x) = x\eu^{x^2}$ em $\R$ tal que
$F(0) = 1$.
\end{exercise}

\begin{exercise}[$\star$]\label{exo:g12:integ:3}
Calcule o \omterm{def:g12:integ:mean}{valor médio} de $f(t) = \sin t$ em $\intcc{0}{\pi}$ e interprete o
resultado em um \omterm{def:g10:functions:graph}{gráfico}.
\end{exercise}

\begin{exercise}[$\star\star$]\label{exo:g12:integ:4}
Usando \omterm{thm:g12:integ:ibp}{integração por partes}, calcule
\[
\int_1^{\eu} \ln t\,\dd t
\qquad\text{e}\qquad
\int_0^{\pi} t \sin t\,\dd t .
\]
\end{exercise}

\begin{exercise}[$\star\star$]\label{exo:g12:integ:5}
Calcule a área da região entre a \omterm{def:g11:quad:parabola}{parábola} $y = x^2$ e a reta
$y = x + 2$.
\end{exercise}

\begin{exercise}[$\star\star$]\label{exo:g12:integ:6}
Seja $I = \displaystyle\int_0^1 \frac{\dd t}{1 + t}$.
\begin{enumerate}
  \item Calcule $I$.
  \item Para $n \in \N$, seja
        $I_n = \displaystyle\int_0^1 \frac{t^n}{1+t}\dd t$. Mostre que
        $I_n + I_{n+1} = \dfrac{1}{n+1}$ e que
        $0 \leq I_n \leq \dfrac{1}{n+1}$.
  \item Deduza que
        $1 - \frac12 + \frac13 - \dots + \frac{(-1)^{n-1}}{n}
        \xrightarrow[n \to +\infty]{} \ln 2$.
\end{enumerate}
\end{exercise}

\begin{exercise}[$\star\star$]\label{exo:g12:integ:7}
A velocidade de um trem (em m/s) durante os primeiros $100$ segundos após
a partida é modelada por $v(t) = 30\bigl(1 - \eu^{-t/50}\bigr)$. Calcule a
distância percorrida nesses $100$ segundos e a velocidade \omterm{def:g11:stat:mean}{média} do trem no
\omterm{def:g10:numbers:interval}{intervalo}.
\end{exercise}

\begin{exercise}[$\star\star\star$]\label{exo:g12:integ:8}
Para $n \geq 1$, seja
$S_n = \dfrac1n \displaystyle\sum_{k=1}^{n} \frac{1}{1 + k/n}$.
\begin{enumerate}
  \item Interprete $S_n$ como uma área de retângulos que aproxima uma
        região sob a curva de $t \mapsto \frac{1}{1+t}$ em
        $\intcc{0}{1}$.
  \item Usando a monotonicidade de $t \mapsto \frac{1}{1+t}$, mostre que
        \[
        S_n \leq \int_0^1 \frac{\dd t}{1+t} \leq S_n + \frac1n\left(1 - \frac12\right),
        \]
        e deduza $\lim\limits_{n\to+\infty} S_n = \ln 2$.
\end{enumerate}
\end{exercise}

\begin{exercise}[$\star\star\star$]\label{exo:g12:integ:9}
\emph{(\omterm{def:g12:integ:area}{Integrais} de Wallis.)} Para $n \in \N$, seja
$W_n = \displaystyle\int_0^{\pi/2} \sin^n t\,\dd t$.
\begin{enumerate}
  \item Calcule $W_0$ e $W_1$.
  \item Escrevendo $\sin^{n+2}t = \sin t \cdot \sin^{n+1} t$ e integrando
        por partes, mostre que $W_{n+2} = \dfrac{n+1}{n+2}\,W_n$.
  \item Deduza $W_2$, $W_3$, $W_4$ e mostre que $(W_n)$ é \omterm{def:g10:functions:variations}{decrescente} e
        positiva.
\end{enumerate}
\end{exercise}

\section{Problema: Arquimedes contra a máquina}

\begin{problem}[{Problema de fim de semana --- a área sob a parábola,
calculada de três maneiras ao longo de vinte e dois séculos, e a
identidade secreta do logaritmo como área}]
\label{pb:g12:integ:1}
Por volta de 240 a.C., Arquimedes calculou a área exata de um
segmento parabólico --- sem \omterm{def:g10:coordgeom:system}{coordenadas}, sem limites, sem álgebra.
Dezenove séculos depois, os retângulos de Riemann refizeram a
conta por puro cerco; e o teorema fundamental do cálculo
(\cref{thm:g12:integ:ftc}) hoje a faz em uma linha. Este problema
disputa as três partidas e depois usa a mesma máquina para revelar
o que o \omterm{def:g12:exp:ln}{logaritmo} realmente é: uma área com uma simetria de
escala.

\medskip
\noindent\textbf{Parte I --- Fluência.}
\begin{enumerate}
  \item Calcule $\displaystyle\int_0^1 (3x^2 - 2x + 1)\,\dd x$,
        \ $\displaystyle\int_1^{\eu} \frac{\dd x}{x}$ \ e
        $\displaystyle\int_0^{\pi/2} \cos x\,\dd x$.
  \item Identifique as formas $u'u$
        (\cref{met:g12:integ:uprime}):
        $\displaystyle\int_0^1 x\,\eu^{x^2}\dd x$ e
        $\displaystyle\int_0^1 \frac{2x}{x^2 + 1}\,\dd x$.
  \item Por partes (\cref{thm:g12:integ:ibp}):
        $\displaystyle\int_0^1 x\,\eu^x \dd x$ e
        $\displaystyle\int_1^{\eu} \ln x\,\dd x$.
  \item Calcule o \omterm{def:g12:integ:mean}{valor médio} (\cref{def:g12:integ:mean}) de
        $\sin$ em $\intcc{0}{\pi}$ --- e note que ele \emph{não}
        é $\frac12$.
  \item Calcule a área entre a reta $y = x$ e a \omterm{def:g11:quad:parabola}{parábola}
        $y = x^2$ em $\intcc{0}{1}$
        (\cref{met:g12:integ:areabetween}).
\end{enumerate}

\medskip
\noindent\textbf{Parte II --- A \omterm{def:g11:quad:parabola}{parábola}, de três maneiras.}
\begin{enumerate}[resume]
  \item A montagem de Riemann para a área sob $y = x^2$ em
        $\intcc{0}{1}$: escreva a soma inferior $L_n$ e a soma
        superior $U_n$ sobre $n$ retângulos iguais (como no
        \cref{exo:g12:integ:8}).
  \item Demonstre por indução (a máquina do
        \cref{pb:g12:seq:1}) a fórmula da soma dos quadrados
        \[
        1^2 + 2^2 + \dots + n^2 = \frac{n(n+1)(2n+1)}{6} .
        \]
  \item Deduza formas fechadas para $U_n$ e $L_n$, calcule seu
        limite comum e conclua: a área é $\frac13$.
  \item Agora a máquina: calcule $\int_0^1 x^2\,\dd x$ pelo
        teorema fundamental, em uma linha. Compare os esforços.
  \item Arquimedes enunciou de outro modo: \emph{um segmento
        parabólico vale $\frac43$ do triângulo nele inscrito}.
        Para o segmento recortado de $y = x^2$ pela corda que
        liga $(-1, 1)$ a $(1, 1)$: calcule a área do segmento por
        uma \omterm{def:g12:integ:area}{integral}, a área do triângulo inscrito (com vértice
        no ponto $(0, 0)$) e verifique a razão do mestre.
  \item O método do próprio Arquimedes: preencha o segmento com o
        triângulo grande $T$, depois com dois triângulos que
        somam $\frac T4$, depois com quatro que somam
        $\frac{T}{16}$, e assim por diante. Some a série
        geométrica e recupere $\frac43 T$ --- a barra de
        chocolate mordida sem fim, do volume do ensino
        fundamental, comida por um geômetra grego.
  \item Uma frase para cada: exaustão (Arquimedes), cerco
        (Riemann), antiderivação (Newton e Leibniz) --- do que
        cada uma precisa e o que cada uma dá?
\end{enumerate}

\medskip
\noindent\textbf{Parte III --- O \omterm{def:g12:exp:ln}{logaritmo} é uma área.}
Para $x > 0$, ponha $A(x) = \displaystyle\int_1^x \frac{\dd t}{t}$.
\begin{enumerate}[resume]
  \item Dê $A(1)$ e $A'(x)$ (\cref{thm:g12:integ:ftc}) e conclua
        que $A$ é exatamente o \omterm{def:g12:exp:ln}{logaritmo natural} da
        \cref{def:g12:exp:ln}.
  \item O milagre da escala: fixe $a > 0$ e estude
        $g(x) = A(ax) - A(x)$. Calcule $g'$ (regra da cadeia),
        deduza que $g$ é constante, avalie a constante --- e
        conclua a \omterm{def:g10:algebra:equation}{equação} funcional
        \[
        \ln(ab) = \ln a + \ln b ,
        \]
        demonstrada por puro cálculo: a área de $1$ até $ab$
        divide-se em cópias reescaladas.
  \item Deduza da questão 14: $\ln(a^n) = n\ln a$ e
        $\ln\frac1a = -\ln a$.
  \item O \cref{exo:g12:integ:8} lê
        $\frac{1}{n+1} + \frac{1}{n+2} + \dots + \frac{1}{2n}$
        como retângulos sob $\frac{1}{1 + t}$: calcule essa soma
        para $n = 10$ (três casas decimais) e compare com
        $\ln 2$. Que somas de sabor harmônico, divergentes termo
        a termo, convergem aqui para uma área?
\end{enumerate}

\medskip
\noindent\textbf{Parte IV --- Acumulação.}
\begin{enumerate}[resume]
  \item Um carro acelera com velocidade $v(t) = 3t^2$ m/s para
        $t \in \intcc{0}{10}$. Calcule a distância percorrida, a
        velocidade \omterm{def:g11:stat:mean}{média} e o instante em que a velocidade
        instantânea iguala a \omterm{def:g11:stat:mean}{média}.
  \item Uma barra de $4$ metros tem densidade linear
        $\rho(x) = 2 + x$ kg/m. Calcule sua massa total e seu
        centro de massa
        $\dfrac{\int_0^4 x\,\rho(x)\,\dd x}{\int_0^4
        \rho(x)\,\dd x}$ --- o ponto de equilíbrio do triângulo
        de papelão do volume do ensino fundamental, finalmente
        calculado com pesos que variam.
  \item As \omterm{def:g12:integ:area}{integrais} de Wallis (\cref{exo:g12:integ:9}): calcule
        $W_0$, $W_1$ e $W_2 = \int_0^{\pi/2} \sin^2 t\,\dd t$
        usando a linearização do \cref{pb:g12:trigo:1}. (Sua
        escada infinita sobe, nos volumes de graduação, até uma
        fórmula de produto para $\pi$ e até a normalização da
        curva em sino.)
  \item Final --- as três faces da integração: uma \emph{área}
        por definição, uma \emph{acumulação} no mundo (distância,
        massa), uma \emph{antiderivada} pelo teorema fundamental;
        e sua história em três nomes. Encerre com a pérola do
        capítulo: que botão corriqueiro da calculadora é, em
        segredo, a área sob $\frac1t$ --- e que identidade essa
        área demonstrou?
\end{enumerate}
\end{problem}
